\fancyhead[L]{\small SCHULZE ÜBER DAS WELTSYSTEM.}
\SourceLeaf{538}
\begin{center}\small Göttingische gelehrte Anzeigen. Stück 21. S. 305..306. 23. Februar 1811.\\[3pt]\rule{.72\linewidth}{.3pt}\end{center}

\emph{Darstellung des Weltsystems. Ein Leitfaden für den Astronomieunterricht an Schulen. Verfasst und mit den wichtigsten bibliografischen Anmerkungen und Nachweisen versehen, um das weitere selbstständige Studium der Astronomie zu erleichtern, von} G. L. \textsc{Schulze}, Prediger in Polenz und Ammelshayn bei Leipzig. Mit 4 Kupfertafeln. 1811. 390 Seiten im Oktavformat. Leipzig, Baumgärtner.

Auch wenn tiefere Kenntnisse der Astronomie ihrer Natur nach nur wenigen Eingeweihten vorbehalten sein können, enthält diese Wissenschaft doch eine solche Fülle höchst interessanter Wahrheiten, die jedem gebildeten Menschen zugänglich sind, dass jeder darin sein ganzes Leben lang eine Quelle vielfältiger Freuden und erhebender Einsichten finden kann. Wer früh Gelegenheit hatte, sich den Weg dorthin zeigen zu lassen, muss das stets dankbar anerkennen. Ein sinnvoll gestalteter Astronomieunterricht für Jugendliche ist daher in dieser Hinsicht sehr wertvoll. Es besteht kein Zweifel, dass sich bei einer stärkeren Verbreitung dieses Unterrichts, als sie bisher bei uns üblich war, sogar manches Talent wecken und zur späteren Weiterbildung anregen ließe. Das vorliegende Lehrbuch ist unseres Erachtens für diesen Zweck sehr gut aufgebaut. Es findet einen glücklichen Mittelweg zwischen den beiden Klippen, an denen man sonst so leicht scheitert, der abschreckenden Trockenheit übermäßiger Abstraktion und der seichten Oberflächlichkeit, die nur Halbwissen fördert und das eigene Denken daran hindert, sich aufzuschwingen. Besonderes Lob verdienen die Hinweise auf die besten astronomischen Schriften zu jedem behandelten Thema. Dort kann jeder, der mehr erfahren möchte, seine Wissbegier stillen. Der Zweck unserer Zeitschrift erlaubt es nicht, bei einem Werk dieser Art weiter ins Einzelne zu gehen. Aber wir müssen gestehen, dass es für uns eine Freude war, bei einem Mann, dessen berufliche Aufgaben ganz anderer Art sind, so gründliche Kenntnisse der Astronomie und eine so große Vertrautheit mit den besten astronomischen Schriften anzutreffen.
